\section{Activities, connections and coefficient norms}
\label{sec:activities}

We specify the local activity space used in \cref{thm:rgnl-current-step}.
These are the conventions of \cite[\S2]{NormalizedContraction}.

\subsection{Chamber sites, pairs and connection domains}

Write \(O_{\mathrm{cub}}\) for the signed permutation matrices on
\(\R^3\). A chamber site is \(x=(c,g)\), with photon projection
\(\pi x=c\) and \(g\in O_{\mathrm{cub}}\). Integration over chamber
sites uses \(\dd x=(\eta_k^3/48)\sum_{c,g}\). We take
\(M\)-cubes centered at \(M\Z^3/\ell\Z^3\); their lattice sites
form a partition. Blocks of every odd side used below are centered
in the same way.

For the three internal neighbors let \(s_0\) change the sign of
the third coordinate, let \(s_1\) exchange the second and third
coordinates, and let \(s_2\) exchange the first and second.
Let \(r_{e_1}\) change the sign of the first coordinate. Set
\[
 x^{(i)}=(c,gs_i)\quad(0\le i\le2),\qquad
 x^{(3)}=(c+\eta_kge_1,gr_{e_1}).
\]
Each neighbor map is an involution. An oriented bond field
\(B\) is antisymmetric under reversal and periodic on the photon
lattice. Write \(B_\mu(c)=B(c,c+\eta_ke_\mu)\), and use the
forward differences
\[
 \partial_\nu B_\mu(c)=
       \eta_k^{-1}\bigl(B_\mu(c+\eta_ke_\nu)-B_\mu(c)\bigr),
 \qquad
 (d\lambda)_\mu(c)=
       \eta_k^{-1}\bigl(\lambda(c+\eta_ke_\mu)-\lambda(c)\bigr).
\]
Thus \((dB)_{\nu\mu}=\partial_\nu B_\mu-\partial_\mu B_\nu\).
Distances between photon sites are sup distances, using the
shortest lift when a distance bound is imposed. For \(0<\alpha<1\)
the three norms are
\[
 \|B\|_1=\|B\|_\infty,\qquad
 \|B\|_2=\|\partial B\|_\infty,\qquad
 \|B\|_3=\|\delta_\alpha\partial B\|_\infty
 =\sup_{\substack{\mu,\nu,c,c'\\0<d(c,c')\le1}}
 \frac{|\partial_\nu B_\mu(c')-\partial_\nu B_\mu(c)|}
      {d(c,c')^\alpha}.
\]
Local versions use only the indicated cells in the given region.
Let \(\square^\dagger\) be the concentric enlargement of an
\(M\)-cube to side \((2L+1)M\). Put
\(\theta_i=e_k^{-3/4+i\varepsilon}\).
The connection domains are
\begin{equation}\label{eq:connection-domains}
\begin{split}
 R_k&=\{B:\|B\|_i<\theta_i,\ 1\le i\le3\},\\
 \widetilde R_k&=\{B_0+B_1:B_1\in R_k,\ B_0\text{ real, and on each }
          \square^\dagger,\\
 &\hspace{25mm} B_0=B_\square^{\mathrm b}+d\lambda_\square
   \text{ with }B_\square^{\mathrm b},\lambda_\square\text{ real and }
          \|B_\square^{\mathrm b}\|_i<\theta_i\},\\
 \mathfrak X_k&=\{C:\|C\|_\infty<e_k^{-1+2\varepsilon},\
                       \|\partial C\|_\infty<e_k^{-1+3\varepsilon}\}.
\end{split}
\end{equation}
Only existence of these representatives is required.
For \(0<c\le1\), the notation \(c\widetilde R_k\) multiplies
all three bounds for both summands by \(c\).
For \(\widetilde R_k^-\) the local real bounds remain
\(\theta_i\), while the global summand has bounds
\((1-e_k^\varepsilon)\theta_i\).

For a short ordered pair \(x=(c,g),x'=(c',g')\), let
\(\tau B(c,c')\) be the arithmetic mean of the sums of \(B\)
along the six coordinate-order paths from \(c\) to \(c'\).
The bond values in this sum are unscaled; define
\[
 U_B(x,x')=\exp\{ie_k\eta_k\tau B(c,c')\},\qquad
 (\delta_{\alpha,B}\psi)(x,x')
       =\frac{U_B(x,x')\psi(x')-\psi(x)}
                    {\max(|c-c'|_\infty,\eta_k)^\alpha}.
\]
For pairs within one cluster the transport is one. Barred fields
use \(U_B^{-1}\). These formulas specify evaluation; the pair
generators in the algebra remain independent.

We use Pauli matrices and
\[
 \gamma_\mu=\sigma_3\otimes\sigma_\mu,\quad
 \gamma_5=\sigma_1\otimes1_2,\quad \Sigma_3=\sigma_3\otimes1_2,
 \qquad
 \mathsf C=i\gamma_2,\quad
 \varkappa=\gamma_5\mathsf C=\sigma_2\otimes\sigma_2 .
\]
Here \(\mathsf C\) is real with
\(\mathsf C^T=\mathsf C^{-1}=-\mathsf C\), while
\(\varkappa\) is a real symmetric involution. The spin group below
is the group generated by \(\gamma_1,\gamma_2,\gamma_3,\gamma_5\).
Its diagonal coordinate action \(r(S)\) is determined by
\(S\gamma_\mu S^{-1}=r(S)_{\mu\mu}\gamma_\mu\).
Connections transform by pullback on oriented bonds.

\subsection{The local Grassmann algebra}
\label{subsec:rg-activity-spaces}

Fix a level $k$, mesh $\eta_k=L^{-k}$ and an admissible side
$\ell$, with $\ell/(LM)$ odd and at least $13$.
Let $\mathcal P_k$ be the collection of face-connected polymers
of $M$-cubes on this torus. A polymer is small when
$d_M(X)\le L$. A small polymer has an injective lift, unique up
to a period translation, by the diameter and lifting estimates
of Appendix~\extref{B}{NormalizedContraction}.
Write $\operatorname{Vol}(X)=M^3|X|_M$.

The independent Grassmann generators are $\psi$, $\bar\psi$ on
the chamber sites $x=(c,g)$ and $\chi$, $\bar\chi$ on ordered
pairs. For distinct chamber sites on the covering lattice, put
$d_X(x,x')=\max(|c-c'|_\infty,\eta_k)$ and restrict pairs to
$d_X\le M$. Pairs on the torus are simultaneous period classes;
their shortest lifts are unique since $2M<\ell$. A pair is
admissible for $X$ when its photon sites and all six rectilinear
paths between them lie in $X$, or when it is $(x,x^{(3)})$ and
one of its photon sites belongs to $X$.
Thus a pair may leave $X$ only along a bond of type $3$.
Pairs inside a cluster have $d_X=\eta_k$ and are admissible when
their common photon site belongs to $X$.
Ordinary generators acquire the sign $(-1)^{n_1+n_2+n_3}$
under translation by $\ell n$. Pair generators acquire the same
sign when both endpoints are translated by $\ell n$.
The kernels are antiperiodic in each ordinary argument and in
each simultaneous pair argument with these conventions.

Let $\xi=(x,\beta,\omega)$ index an ordinary leg and
$\zeta=(x,x',\beta,\omega)$ a H\"older leg, where
$\omega=0$ is unbarred and $\omega=1$ is barred.
We use the index measures
\[
 \nu(\xi)=\eta_k^3/48,\qquad
 \nu(\zeta)=(\eta_k^3/48)^2.
\]
Integrals sum the spin and bar indices as well. An activity
has a unique expansion with separately antisymmetric kernels,
\begin{equation}\label{eq:rg-activity-kernels}
 \mathscr E(X;\mathscr A,\mathscr C)
 =\sum_{n,m}\frac1{n!m!}\int\dd\xi\,\dd\zeta\,
      \mathscr E_{nm}(X;\mathscr A,\mathscr C;\xi;\zeta)
                       \psi(\xi_1)\cdots\psi(\xi_n)
                       \chi(\zeta_1)\cdots\chi(\zeta_m).
\end{equation}
Barred generators are included in this notation.
The generator weights, as distinct from the index measures, are
\[
 h_{k,1}=e_k^{-1/4},\qquad h_{k,2}=e_k^{-1/4+\varepsilon}.
\]
For an element $F$ the weighted coefficient norm is
\begin{equation}\label{eq:rg-activity-coefficient-norm}
 \|F\|_{h_k}
       =\sum_{n,m}\frac{h_{k,1}^nh_{k,2}^m}{n!m!}
                                   \int\dd\xi\,\dd\zeta\,|F_{nm}|.
\end{equation}
In particular a generator has norm $h_{k,i}$;
the weight $\nu$ enters only when coefficients are written as
integral kernels.

An activity family in $K_k^{x,\C}$ is balanced: every monomial
has as many barred as unbarred legs. Its kernels are holomorphic
on $\widetilde R_k\times\mathfrak X_k$, vanish unless all ordinary
legs lie over $X$ and all pairs are admissible for $X$, and read
the two connections only on bonds meeting $X$. The norm is
\begin{equation}\label{eq:rg-activity-norm}
 \|\mathscr E\|_k^x
  =\sup_{X\in\mathcal P_k}\exp(\kappa d_M(X))
    \sup_{(\mathscr A,\mathscr C)\in\widetilde R_k\times\mathfrak X_k}
                          \|\mathscr E(X;\mathscr A,\mathscr C)\|_{h_k}.
\end{equation}
We require $\kappa\ge\kappa_0$ and finite norm.

The required symmetries act on both the connections and the
generators. For $\tilde g=(S,a)$ with the coordinate sign action
$r(S)$ and $a\in M\Z^3$, put
\[
 \tilde g(c,g)=(r(S)c+a,r(S)g),\qquad
 a_{\tilde g}\psi(x)=S\psi(\tilde g^{-1}x),\qquad
 a_{\tilde g}\bar\psi(x)=(S^{-1})^T\bar\psi(\tilde g^{-1}x).
\]
The same spin matrices act on pair generators after transforming
both endpoints. Local gauge transformations multiply an unbarred
ordinary or pair leg by $\exp(ie_k\Lambda(c))$, using the photon
site of its first point, and a barred leg by the inverse phase.
Write this automorphism as $G_\Lambda$.
Charge conjugation and the antilinear reality map act by
\[
 \mathrm{Cg}\,\psi=\mathsf C\bar\psi,\qquad
 \mathrm{Cg}\,\bar\psi=-\mathsf C\psi,\qquad
 \mathbf K\psi=\varkappa\psi,\qquad
 \mathbf K\bar\psi=\varkappa\bar\psi,
\]
and by the same rules on the H\"older generators.

The covariance requirements on a family are
\begin{equation}\label{eq:rg-activity-symmetries}
 \begin{split}
 a_{\tilde g}\mathscr E(\tilde gX;\tilde g\mathscr A,
                                      \tilde g\mathscr C)
                                  &=\mathscr E(X;\mathscr A,\mathscr C),\\
 \mathrm{Cg}\,\mathscr E(X;-\mathscr A,-\mathscr C)
                                  &=\mathscr E(X;\mathscr A,\mathscr C).
 \end{split}
\end{equation}
For real scalars $\lambda,\lambda'$ and two pairs in the domain
which satisfy $\mathscr A^1=\mathscr A-d\lambda$ and
$\mathscr C^1=\mathscr C-d\lambda'$ on bonds at distance at
most one from $X$, the local gauge requirement is
\[
 G_{\lambda+\lambda'}\mathscr E(X;\mathscr A^1,\mathscr C^1)
                               =\mathscr E(X;\mathscr A,\mathscr C).
\]
The real subspace $K_k^x$ has in addition
\[
 \mathbf K\mathscr E(X;-\overline{\mathscr A},-\overline{\mathscr C})
                               =\mathscr E(X;\mathscr A,\mathscr C).
\]
All requirements refer to the generators before evaluating fields.
They are essential to the form of the counterterms.

The spaces $K_k^\C,K_k$ without the second connection argument
have the same definitions on $\widetilde R_k$. Restriction
$\mathscr E|_0(X;\mathscr A)=\mathscr E(X;\mathscr A,0)$ maps
$K_k^{x,\C}$ into $K_k^\C$ and has norm at most one. The real
subspaces also restrict. Analyticity, locality and the symmetries
follow by setting the second argument and its gauge function to
zero; the norm bound follows from the smaller supremum.

\subsection{Total connections}
\label{subsec:rg-total-connection-domain}

The comparison domain $\widetilde G_k$ consists of fields $B_0+B_1$
with $B_0$ real and $\|B_1\|_\infty<e_k^{-1+\varepsilon}$,
where on each $\square^{\sim5}$ the field $B_0$ has a real-gauge
representative of supremum norm below $e_k^{-1+\varepsilon}$.
Here $\square^{\sim5}$ has the same center and side $11M$.
The real-gauge representatives in the domains defined above may
be chosen independently on each enlarged cube.

Assume $L\ge5$, $0<\varepsilon\le1/200$ and
\begin{equation}\label{eq:rg-domain-charge-conditions}
 e_k^\varepsilon\le\tfrac14,\qquad
 e_k^{1/4}+e_k^\varepsilon\le1,\qquad
 4e_k^{1/4+2\varepsilon}+2e_k^{3\varepsilon}\le1.
\end{equation}
Then
\begin{equation}\label{eq:rg-total-domain-inclusion}
 \widetilde R_k+\mathfrak X_k\subset\widetilde G_k.
\end{equation}
Indeed, write $\mathscr A=A_0+A_1$ with $A_1\in R_k$ and
$A_0$ locally real-gauge equivalent on $\square^\dagger$ to
$A_\square^b$ with the three threshold bounds. Restrict the
representative to $\square^{\sim5}\subset\square^\dagger$.
Its supremum norm is below
$e_k^{-3/4+\varepsilon}\le e_k^{-1+\varepsilon}$, and
\[
 \|A_1+\mathscr C\|_\infty
 <e_k^{-3/4+\varepsilon}+e_k^{-1+2\varepsilon}
 =(e_k^{1/4}+e_k^\varepsilon)e_k^{-1+\varepsilon}
 \le e_k^{-1+\varepsilon}.
\]
This is the required decomposition of the total connection.
Real gauge transformations can absorb any real constant part,
locally on the nonwrapping cubes.

Put $\mathfrak i(B)=e_k\|\operatorname{Im}B\|_\infty$ and
$q_d(B)=e_k\|dB\|_\infty$. The same decomposition gives
\begin{equation}\label{eq:rg-total-imaginary-curvature-bounds}
 \mathfrak i(\mathscr A+\mathscr C)
       \le e_k^{1/4+\varepsilon}+e_k^{2\varepsilon}\le1,
 \qquad
 q_d(\mathscr A+\mathscr C)
       \le4e_k^{1/4+2\varepsilon}+2e_k^{3\varepsilon}\le1.
\end{equation}
For the curvature, every plaquette lies in an enlarged cube,
and $dA_0=dA_\square^b$ there. This costs at most
$2\|\partial A_\square^b\|_\infty$; the global small field
costs another $2\|\partial A_1\|_\infty$, and the external
field costs $2\|\partial\mathscr C\|_\infty$.
The imaginary bound uses only that $A_0$ is real.

The domains are open: the strict inequalities define open
sets for the global summand, which can absorb a sufficiently
small complex perturbation while the local real representative
is kept fixed. The norm-defined sets are invariant under sign
and conjugation; the local real-gauge definitions are also
invariant under real gauges and real constant additions.
Coordinate sign changes and $M$ translations preserve them.
Their derivative and H\"older norms are preserved, with a
possible reversal of a bond and translation of its base point.
Thus \eqref{eq:rg-total-domain-inclusion} applies also at the
signs, conjugates and real gauges used below.



\subsection{The operator and photon data}

The injection and cluster mean are
$(Jv)(c,g)=v(c)$ and $(J^*f)(c)=48^{-1}\sum_gf(c,g)$.
The normalized average $\mathcal Q_k=Q_kJ^*$, its reverse
$\mathcal Q_k^\sharp=\mathcal Q_k(-B)^T$, the one-step average
$Q$, and the chamber operator $D_B$ are given in
\cite[\S1]{FermionBlockManuscript}. Put $b_k=(1-L^{-1})/(1-L^{-k})$ for $k\ge1$. At these levels set
\[
 S_k(B)=(D_B+\mu+b_k\mathcal Q_k^\sharp\mathcal Q_k)^{-1},\quad
 H_k(B)=b_kS_k(B)\mathcal Q_k^\sharp(B),\quad
 D_k(B)=b_k-b_k^2\mathcal Q_k(B)S_k(B)\mathcal Q_k^\sharp(B).
\]
The one-step projection is $P_L(B)=Q^T(-B)Q(B)$. The barred map is $\bar H_k(B)=\mathsf C^{-1}H_k(-B)\mathsf C$.
At $k=0$, $H_0=\bar H_0=1$ and $D_0(B)=D_B+\mu$.
The symbol $\Sspace(a,\ell)$ denotes the four-component spinor
space on the $a$-spaced torus, with measure $a^3$; the superscript
$\mathrm{ch}$ adds chamber labels and divides the measure by $48$.
All spaces have the antiperiodic convention stated above.
For later use, the hopping notation of \cite[\S2.1]{FermionBlockManuscript} is
\[
 R_i(B)f(x)=U_B(x,x^{(i)})f(x^{(i)}),\qquad
 \mathcal L_B=\sum_{i=0}^3w_i(1-R_i(B)),\qquad
 M_i(c,g)=\tfrac{w_i}{2}(\rho+\gamma\cdot g\mathbf n_i),
\]
where $(w_0,w_1,w_2,w_3)=(6,6\sqrt2,6\sqrt2,6)$ and
$(\mathbf n_0,\mathbf n_1,\mathbf n_2,\mathbf n_3)=(e_3,(e_2-e_3)/\sqrt2,(e_1-e_2)/\sqrt2,-e_1)$.
Thus $\widehat d_i(c,g)=-g\mathbf n_i$. On a nearest-neighbor photon bond
write $u_B(c,c')=\exp(ie_k\eta_kB(c,c'))$; this is the transport
$U_B(x,x^{(3)})$ when $c,c'$ are its endpoints.

For photon fields use the Landau reconstruction $H_k^{\mathrm{ph}}$,
the constrained minimum $A_k^{\min}$ and the real coordinate map
$C:V_k\to\mathcal Z$ of
\cite[Theorem~A]{MaxwellBlockManuscript}. Write
$C_k^{1/2}=\widetilde\Delta_k^{-1/2}$ and
$\mathrm{Sq}_k=CC_k^{1/2}$. The coordinate basis $\{\Upsilon\}$ is
orthonormal and has the specified block or face-bond supports.
The law $\mu_I$ is the product standard Gaussian in these coordinates.
Its restriction to $\mathrm{Box}_k$ has probability
$\exp(-\ell^3\delta\varepsilon_k^0)$, with
$0\le\delta\varepsilon_k^0\le12\exp(-p_{0,k}^2/2)$.
The symbols $C$ and $\mathscr C$ always denote, respectively,
this real coordinate map and the external connection.

The common constant $c_\gamma$ bounds both
\[
 \sup_\Upsilon\sum_{\Upsilon'}\exp[-\gamma_wd(\Upsilon,\Upsilon')],
 \qquad \sup_b\sum_{b'}\exp[-\gamma_wd(b,b')],
\]
where the second sum is over unit bonds. The constants $C_w$ and
$\gamma_w$ dominate the photon bounds and the corresponding
fermion bounds as stipulated in hypothesis~(i).
